\section{Energy and boundary flux}\label{ch:mass:section}

Fix a prepared strict rest datum, including its face collars, offsets and
weights.  All limits in this section hold for that datum and for a fixed
sufficiently small $\epsilon>0$, with $n\to\infty$ only after the
fixed-$n$ exhaustion in Proposition~\ref{ch:ell:neutral-input}.
Constants may depend on the datum and on $\epsilon$, but not on $n$,
unless expressly stated otherwise.

\begin{proposition}[Energy of the comparison metric]\label{ch:mass:comparison-energy}
Let $(f,t)$ be the global solutions of Equation~\eqref{ch:ell:system} supplied
by Proposition~\ref{ch:ell:neutral-input}, and let
$\widetilde g_n=e^{4\epsilon}\hat g_n$ and the comparison fields be as in
Proposition~\ref{ch:cmp:charged-comparison}.  In the normalized end coordinates
$z=e^{2\epsilon}y$, these data satisfy
\[
 \widetilde g_n-\delta=O_2(|z|^{-1}),\qquad
 \widetilde X_{1,n},\widetilde X_{2,n}=O_1(|z|^{-2}),\qquad
 R_{\widetilde g_n}\in L^1(X,dV_{\widetilde g_n}).
\]
Their ADM energies are finite and obey
\begin{equation}\label{ch:mass:energy-bound}
 \widetilde E_n\le e^{2\epsilon}(E_*+\eta_n),
 \qquad \eta_n\ge0,\qquad \eta_n\longrightarrow0.
\end{equation}
No uniform bound in $n$ on the solutions' classical regularity constants
is needed for this conclusion.
\end{proposition}

\begin{proof}
The charged preparation belongs to the reduced neutral class by
Lemma~\ref{ch:ell:hypothesis-check}, and
Proposition~\ref{ch:ell:neutral-input} supplies its exact physical
endpoint after exhaustion at fixed $\epsilon,n$. We may therefore
apply the neutral end and flux results. We verify the relevant
hypotheses before carrying out the additional scaling and field
comparison.

\emph{The end and its ADM flux.}
The prepared tensor $K$ is compactly supported, the end satisfies
$g-\delta=O_j(r_y^{-1})$ and $R_g=O_j(r_y^{-3-\delta})$ for every
fixed $j$, and the neutral offset inequality holds. The compact
height bound is used only on the supports of $K,C$. The profiles,
quadratic form and flux are exactly those identified in
Lemma~\ref{ch:cmp:scalar}, with the capillarity parameter
$\tau_0=\ell^{3/2}$ denoted by $\tau$ in the neutral argument.
Lemmas~\ref{n3:limit:end} and~\ref{n3:limit:geometry} consequently give
\[
 \hat g_n-\delta=O_2(r_y^{-1}),\qquad
 R_{\hat g_n}\in L^1(X,dV_{\hat g_n}),\qquad
 V=4\nabla_gt+O(r_y^{-3}),
\]
and the finite flux identity
\begin{equation}\label{ch:mass:adm-flux-identity}
 E_{\hat g_n}=E_*-\frac{\mathfrak F_n}{8\pi},\qquad
 \mathfrak F_n=\lim_{s\to\infty}
                   \int_{r_y=s}V_{\nu_s}\,dA_g.
\end{equation}
Here $\nu_s$ points toward infinity. These statements require only
fixed-$n$ end estimates. In particular the exponential graph error
has zero ADM contribution, and the leading conformal change
$4t\delta$ has energy flux $-8\partial_i t$, as in the proof of
Lemma~\ref{n3:limit:geometry}.

\emph{The lower flux bound.}
Equation~\eqref{ch:ell:system} has the source
$\delta_0\cT+\rho+\delta_0\tau_0a\sigma-m_0(t)\cP$ and the
boundary flux used in Lemma~\ref{n3:limit:flux}. The two face
identities are $F-P_B=H$, $w_\nu=a$ on black faces and
$F-P_B=-H$, $w_\nu=-a$ on white faces. The signed slope estimates
\eqref{ch:ell:slopes} hold on every existing face; no condition is
needed on an absent color. The weight $\rho$ has a fixed positive
minimum on the compact collars, the penalty $C_n$ is polynomial,
and the weights $\rho_0,\rho_0^2,\rho_1^2$ are integrable because
$\gamma>3/4$. These are precisely the boundary, bulk and penalty
hypotheses of that lemma. It gives
\begin{equation}\label{ch:mass:vanishing-error}
 \mathfrak F_n\ge-\zeta_n,\qquad
 0\le\zeta_n\le C(1+n)^N
         \left(\ell+\frac{\ell^2}{\tau_0}\right)
       =C(1+n)^N(\ell+\ell^{1/2})\longrightarrow0.
\end{equation}
Its proof absorbs the boundary contribution by the polynomial
collar trace before estimating the floor penalty. Thus the constants
in this bound do not involve the possibly much larger fixed-$n$
classical regularity constants.

\emph{Scaling and electromagnetic fields.}
The constant rescaling $\widetilde g_n=e^{4\epsilon}\hat g_n$ is
asymptotically Euclidean in coordinates $z=e^{2\epsilon}y$, and
\begin{equation}\label{ch:mass:scaling}
 \widetilde E_n=e^{2\epsilon}E_{\hat g_n}.
\end{equation}
Indeed the coordinate derivative and sphere area factors in the ADM
integral are $e^{-2\epsilon}$ and $e^{4\epsilon}$, respectively.
This also gives the asserted metric decay and preserves scalar-curvature
integrability.

The original flux forms have components $O_1(r_y^{-2})$ on this end.
The forms in \eqref{ch:cmp:new-flux-forms} are constant linear
combinations of them and have the same decay. Define their vector
fields by $\iota_{\widetilde X_{i,n}}dV_{\widetilde g_n}
=\widetilde\alpha_i$. The metric coefficients are constant at
infinity up to $O_2(r_y^{-1})$, so this definition followed by the
fixed coordinate dilation gives
$\widetilde X_{i,n}=O_1(|z|^{-2})$. Their integrals are unchanged
by this coordinate change, so both total charges remain $Q_E,Q_B$.
Proposition~\ref{ch:cmp:charged-comparison} supplies the charged
scalar-curvature inequality for these very fields and metric.

Combining \eqref{ch:mass:adm-flux-identity}--\eqref{ch:mass:scaling}
proves \eqref{ch:mass:energy-bound} with
$\eta_n=\zeta_n/(8\pi)$. Throughout this argument the datum and
$\epsilon$ are fixed and the exhaustion precedes $n\to\infty$.
\end{proof}
